% Root and D full literal audits passed; independent rational control passed.
\subsection{Five states opposite the three-element right-zero band}

\begin{lemma}[A vertex-fixing scalar map of a double-cover graph]
\label{C:double-cover-graph-scalar}
Let a connected finite graph $A$ have a free cellular $C_2$-action
on vertices and edges. For every odd positive integer $d$, there
is a $C_2$-equivariant self-map fixing every vertex and inducing
$d$ times the identity on $H_1(A;\mathbb Z)$.
\end{lemma}
\begin{proof}
The quotient $B=A/C_2$ is a connected graph, and its connected
double cover is classified by a surjection
$\eta:F_r=\pi_1(B)\to C_2$.
Choose a free basis $a,b_1,\ldots,b_{r-1}$ with
$\eta(a)=1$ and $\eta(b_j)=0$; replacing another odd basis
element by its product with $a^{-1}$ gives such a basis.
The endomorphism $a\mapsto a^d$, $b_j\mapsto b_j^d$
preserves $\eta$. Represent it by a cellular map of $B$
that fixes a chosen maximal tree pointwise. Its lift fixing
a chosen lifted base vertex is equivariant for the deck
involution, and fixes the two lifts of every tree vertex,
hence every vertex of $A$.

A Schreier basis for the kernel of $\eta$ is
\[
 a^2,\qquad b_j,\qquad ab_ja^{-1}.
\]
The first two types are multiplied by $d$ after abelianization.
Writing $d=2k+1$, the image of the last type is
$a^{2k}(ab_j^da^{-1})a^{-2k}$, so its abelianization is
also multiplied by $d$. This is exactly the claimed action
on the cover's first homology.
\end{proof}

\begin{lemma}[Removing constant outside stars]
\label{C:right-zero-constant-star-removal}
Let $R$ be the identity-adjoined right-zero three-band acting
on the left of a finite biset $S$, with core $U$.
Suppose that every outside state on which the three minimal
actors agree has its common image at the same $u\in U$.
Collapse all those states to $u$. This gives an actual right
$M$-equivariant quotient of the natural action graph.
Either $u$ is fixed by all of $M$, or the induced map on
collage cofibers is an isomorphism on second rational
cohomology by pullback.
\end{lemma}
\begin{proof}
Let $A$ be the vector space on these outside states, with
the partial right action that sends a basis state to zero
when it enters the core. Such a state cannot map to a
nonconstant outside state, by commutation with the three
retractions. If it maps to another constant outside state,
the image of $u$ is $u$; if it maps into the core, that
image is also the image of $u$. Thus the quotient relation
is right-invariant.

The graph-chain kernel of the collapse has homology only
in degree one, where it is $A\otimes\widetilde{\mathbb Q^3}$.
The second factor records the three edge colors and is
fixed by every right actor. If $A\otimes_{\mathbb QM}\mathbb Q$
is nonzero, its state-transition relations have a nonempty
component with no transition into the core. Each actor
therefore sends a state of that component to another
constant outside state attached to $u$; commutation forces
that actor to fix $u$. Hence $u$ is a common fixed state.
Otherwise the kernel's first-homology coinvariants are zero.
After the free bar tensor and the cofiber suspension, the
kernel complex has zero homology in both degrees one and
two. Its homology exact sequence therefore makes the actual
cofiber comparison an isomorphism in degree two. Duality
over $\mathbb Q$ proves the stated cohomology isomorphism.
\end{proof}

\begin{theorem}[The remaining four-element corner on five states]
\label{C:right-zero-three-five-small-sources}
Let $R$ be the identity-adjoined three-element right-zero band,
let $|M|\leq9$, and let $|S|\leq5$ in an $(R,M)$-biset.
If its collage has one-dimensional second rational cohomology,
that class has square zero. In particular the collage cannot
have type $\mathbb CP^2$.
\end{theorem}
\begin{proof}
Use the natural graph with core $U$, one additional vertex
for every outside state, and three colored edges from that
state to its three retraction images. Right actors preserve
colors and send an edge at unit speed to another edge or
collapse it. A singleton core gives a common fixed state.
If all three retractions coincide, their common biset
retraction gives an ordinary-act retract and a suspension
quotient. Lemma~\ref{C:nine-unbounded-act-cups}, together
with the assumed rank one in degree two, kills the square.
This also includes the trivial left action.

The proof of Theorem~\ref{C:right-zero-three-four-states}
works with this same rational rank-one hypothesis when
$|S|\leq4$: replace its ordinary-act conclusion in the
common-retraction branch by that universal zero-product
lemma. All its other branches already give zero products
or a contradiction in second rational homology.
It therefore remains to treat exactly five states.

If constant outside states occur only at one core vertex,
Lemma~\ref{C:right-zero-constant-star-removal} either gives
a common fixed state or identifies the second cohomology
with that of a graph on at most four states. Its rank remains
one, so the preceding
paragraph finishes that case. If constant outside states
occur at two different core vertices and a nonconstant
outside state also exists, the state budget forces exactly
two core vertices, one constant outside state at each, and
one nonconstant outside state. We treat this configuration
last. Hence for the moment all outside states may be
assumed nonconstant.

\emph{Two core vertices.}
There are three outside states. Each has a nonconstant
binary word $v_s\in\{0,1\}^3$ recording its colored
endpoints. A vertex-constant actor sends every state to
that constant core vertex, so kills all first homology
of the connected graph and makes the collage's second
homology zero. Without such actors, the core actions are
identity or swap. If no swap occurs, a core vertex is a
common fixed state.
If a swap occurs, the distinct outside words occur in
complementary pairs. There are at most three words, so
there is exactly one complementary pair. Identify all
outside vertices having the same word, preserving each
colored edge. This is an equivariant graph quotient and
is onto on first homology: choosing one representative
for each word gives an underlying graph section.
The quotient graph has four vertices and six edges, and
its right-action image is the free $C_2$ action interchanging
both pairs of vertices. Its orbit graph has two vertices
and three edges, so its first rational homology coinvariants
have dimension two. The original connected graph's
coinvariants surject onto them. The graph-Tor comparison
contradicts one-dimensional second homology.

\emph{Three core vertices.}
There are two nonconstant outside words $v_z,v_w$.
If one word uses all three core vertices, relabel it as
the identity word. If the other also uses all three, a
nonidentity core actor must carry the first word to the
second and vice versa. It is a transposition on the three
core vertices; otherwise its square could not be a word
present in the action. A core vertex is therefore fixed
by every actor. If the two full words agree, all core
actors are already the identity.
If the other word uses only two vertices, the only
possible nonconstant core actors are the identity and
that one rank-two map $\sigma$. A constant core actor
collapses every state, killing second homology.
Closure then forces $\sigma^2=\sigma$; a constant square
would give the same contradiction, and a rank-three square
is impossible. Thus a vertex in its nonempty image is fixed
by every actor. This disposes of every full word.

Suppose both words use two core vertices. If they use the
same pair, the graph has a cyclic component on those two
core vertices and the two outside states, and one isolated
core vertex. An actor collapsing that component's first
homology makes its coinvariants zero, so degree-two classes
pull back from the two-state component act and square to
zero. Otherwise the four-state word analysis shows either
a common fixed core vertex, or complementary words with
the free $C_2$ graph action just described.
In the latter case Lemma~\ref{C:double-cover-graph-scalar}
gives a degree-three scalar map on the component fixing
all its vertices. Extend it by the identity on the isolated
vertex. This is strictly equivariant: actors restrict to
the identity or deck involution on the component, and
send the isolated core vertex only to core vertices.
The component-eigenvalue lemma and the two-state finite-line
lemma kill the rank-one square.

If the endpoint pairs differ, their union is all three
core vertices, and the graph is connected with two
independent cycles, one for each outside state. If the
two words repeat different pairs of color positions, an
actor cannot send one outside state to the other: applying
a vertex map preserves equal color positions. Each cycle
therefore carries its own zero-or-one character, recording
whether that outside state is fixed or enters the core.
Nonzero coinvariants require at least one outside state
to be fixed by all actors. Its core endpoints are then
common fixed states, again killing products.

If the two words repeat the same pair of color positions,
their cycle module is the partial permutation module on
$\{z,w\}$. A nonzero coinvariant class determines a
nonempty transition component that is invariant and never
enters the core. A singleton such component gives a fixed
outside state and hence a fixed core endpoint.
Otherwise both outside states stay outside under every
actor. Regard each as an oriented edge from its repeated
endpoint to its remaining endpoint. The two edges form
an oriented path on the three core vertices. Every actor
sends each edge to an edge with its orientation, without
collapse. For the two possible star orientations the
middle vertex is fixed, and for a directed path only the
identity edge assignment is possible. Thus the middle
vertex is a common fixed state in every orientation.

\emph{Four core vertices.}
There is one nonconstant outside state. Three distinct
endpoints give a forest, whose component quotient is an
ordinary act cone and has zero degree-two products.
Two distinct endpoints give one unicyclic component.
If an actor sends the outside state to the core it kills
the cycle coinvariants; the class then pulls back from
the component act, and Lemma~\ref{C:nine-unbounded-act-cups}
applies. Otherwise that outside state is fixed by every
actor, and either of its endpoints is a common fixed state.
This finishes the case of no constant outside states.

\emph{Two constant stars and one nonconstant state.}
Write the core as $u_0,u_1$, the constant outside states
as $w_0,w_1$, and the remaining state as $z$.
The word at $z$ has two equal entries and one other entry.
A core swap is impossible because its complementary word
is absent. Without both constant core maps, a core vertex
is fixed, so assume both occur.

Put $V=\widetilde{\mathbb Q^3}$, the color-difference space.
The first homology of the connected graph decomposes as
$V_{w_0}\oplus V_{w_1}\oplus L_z$, where $L_z\subset V$
is the one-dimensional line comparing the equal color
positions of $z$. A nonzero invariant functional has
restrictions $\lambda_0,\lambda_1,\lambda_z$ to these
summands. For an actor constant at $u_i$, each $w_j$ or
$z$ is sent either to $u_i$ or $w_i$.
Invariance under the constant actors at both vertices
forces $\lambda_0=\lambda_1=\lambda\ne0$, and forces
both $w_0,w_1$ to be sent to $w_i$ by every such actor.
It also forces every core-fixing actor to fix both $w_j$
and $z$, hence to be the identity on all five states.

Every remaining actor is therefore $c_{i,\delta}$:
it sends the two core states to $u_i$, both constant
outside states to $w_i$, and $z$ to $u_i$ if $\delta=0$
or $w_i$ if $\delta=1$. Invariance says
$\lambda_z=\delta\lambda|_{L_z}$.
If all actors have the same $\delta$, the invariant
functional space has dimension two, since $\lambda\in V^*$
is arbitrary. This contradicts the rank-one assumption.
Otherwise closure under
$c_{i,\delta}c_{j,\epsilon}=c_{j,\delta}$ gives all
four actors and the identity. It forces
$\lambda|_{L_z}=0$ and $\lambda_z=0$; their invariant
functional space is exactly one-dimensional.

Now identify in $R$ the two right-zero labels occupying
the repeated positions of the word at $z$. The left action
on every state factors through this actual quotient
$R\to\{1,p,q\}$. Simultaneously replace $M$ by its
five-element action image. The resulting actual collage
is precisely the thirteen-arrow control of
Lemma~\ref{C:nine-five-state-graph-square}.
On first graph-homology coinvariants its pullback identifies
the target generator with the line
$\lambda_0=\lambda_1=\lambda$, $\lambda|_{L_z}=0$,
$\lambda_z=0$ just computed. The actual functor therefore
induces an isomorphism on second rational cohomology.
The signed rational square primitive of that literal
control pulls back and finishes the last case.
\end{proof}

\begin{corollary}[All order-four opposite corners at the first five-state profiles]
\label{C:four-corners-five-state-eighteen}
Every triangular two-object category with $|N|=4$,
$|S|\leq5$, and $|M|\leq9$ is excluded from type
$\mathbb CP^2$. In particular the seventeen-arrow profile
$(8,4,5)$ is excluded completely.
\end{corollary}
\begin{proof}
Projective trivial right modules are covered by the accepted
finite-target opposite-corner comparison and
Theorem~\ref{C:finite-target-small-act-twenty}.
The five nonprojective corners with a two-element minimal
band are covered by
Theorem~\ref{C:four-corner-two-band-five-states}.
The remaining corner is the one treated above.
\end{proof}
